\documentclass[10pt]{article}
\usepackage[margin=0.7in]{geometry}
\usepackage[T1]{fontenc}
\usepackage{lmodern}
\usepackage{amsmath,amssymb}
\usepackage{enumitem}
\setlist[enumerate]{leftmargin=*, itemsep=7mm}
\setlist[enumerate,2]{label=(\Alph*), leftmargin=2em, itemsep=0.05em, topsep=0.25em}
\pagestyle{plain}

\begin{document}
\begin{center}
  {\large\bfseries Video 5 Quiz: Metric Spaces and Topology}\\[0.25em]
  \textit{Choose the best answer for each question.}
\end{center}
\noindent Name:\ \rule{3.2in}{0.4pt}\hfill Date:\ \rule{1.25in}{0.4pt}
\vspace{0.7em}
\begin{enumerate}
  \item What intermediate location is used in the video’s travel interpretation of the triangle inequality?
  \begin{enumerate}
    \item Where you work.
    \item Where you live.
    \item Where you buy groceries.
    \item A nearby gas station.
  \end{enumerate}

  \item Which sequence represents the video’s example in $\mathbb{R}$ with the absolute-distance metric?
  \begin{enumerate}
    \item $x_n=(1+\tfrac{1}{n})^n$.
    \item $x_n=\tfrac{1}{n!}$.
    \item $x_n=(-1)^n$.
    \item $x_n=n^2$.
  \end{enumerate}

  \item In the counterexample with $X=\mathbb{Q}$, why does the Cauchy sequence not converge in that metric space?
  \begin{enumerate}
    \item Its terms eventually stop being rational.
    \item Its limit is not an element of $\mathbb{Q}$.
    \item The absolute-distance metric is not symmetric on $\mathbb{Q}$.
    \item A Cauchy sequence can never have a limit.
  \end{enumerate}

  \item What does the statement $P$ on the open-ball slide say students will be asked to prove in the homework?
  \begin{enumerate}
    \item Every point in an open ball is inside a smaller open ball contained in the original ball.
    \item Every open ball is a closed subset of the metric space.
    \item Every convergent sequence stays inside one open ball.
    \item Every metric space has only finitely many open balls.
  \end{enumerate}
\end{enumerate}
\end{document}
